\chapter{Procesamiento distribuido: transformación, resolución de entidades e indicadores}
\label{cap:procesamiento}

El presente capítulo describe la parte computacionalmente más densa del
sistema. Se recorren la validación y normalización de bronze a silver, la
resolución de identidad entre un mismo investigador presente en ORCID y en
un CVN sintético, y el cálculo de indicadores de investigación publicados
en gold. Es el capítulo donde más se demuestra el procesamiento distribuido
y donde el enlace por identificador ORCID sembrado en el
Capítulo~\ref{cap:ingesta} rinde su fruto, al convertirse en la evidencia
que permite fusionar los registros de una misma persona.

\section{Arquitectura de la transformación}
\label{sec:procesamiento_arquitectura}

La transformación se implementa como tres trabajos de Spark, encadenados
en un único flujo de Airflow llamado \texttt{transform\_publish}. El
primero, \texttt{bronze\_to\_silver}, valida y normaliza bronze a silver.
El segundo, \texttt{silver\_to\_gold}, calcula los indicadores de
investigación. El tercero, \texttt{publish\_gold\_to\_postgres}, publica
gold en una base de datos relacional. Cada uno corre en su propio pod,
igual que la ejecución de Spark verificada en el
Capítulo~\ref{cap:infraestructura} y la ingesta del
Capítulo~\ref{cap:ingesta}. Una restricción nueva condiciona todo el
código de este capítulo. La distribución de Spark empleada trae su propio
intérprete de Python, en la versión 3.10, distinta de la versión 3.14 que
exige el resto del proyecto. El código de este capítulo debe, por tanto,
ser compatible con Python 3.10, comprobado mediante una prueba dedicada
que analiza la sintaxis de cada módulo contra esa versión antes de
ejecutarlo.

Las dos primeras etapas escriben directamente tablas Iceberg bajo los
espacios de nombres \texttt{lakehouse.silver} y \texttt{lakehouse.gold},
sin aterrizaje intermedio de ficheros crudos, como quedó establecido en la
Sección~\ref{sec:infraestructura_iceberg}. La tercera etapa publica gold en
una base de datos PostgreSQL dedicada, para que el panel de visualización
del Capítulo~\ref{cap:visualizacion} nunca necesite consultar Iceberg ni
Spark directamente. Las tres etapas comparten una única imagen de Spark,
construida añadiendo a la imagen del Capítulo~\ref{cap:infraestructura}
las dependencias de validación del proyecto y el controlador JDBC de
PostgreSQL. El número de ejecutores, la memoria de cada ejecutor y el
número de particiones de mezcla de cada trabajo son parámetros explícitos
del flujo. El Capítulo~\ref{cap:visualizacion} reutiliza esos parámetros
como variable de la prueba de rendimiento que allí se describe.

\section{Validación y normalización de bronze a silver}
\label{sec:procesamiento_silver}

El primer trabajo lee cada fuente de bronze por separado como líneas de
texto, para preservar el documento original y evitar que la inferencia de
esquema de Spark colapse el contenido heterogéneo de un CVN. Cada línea
se deduplica por identificador de registro, conservando la copia con la
fecha de aterrizaje más reciente, antes de analizar su contenido, así que
ningún trabajador llega a analizar dos veces un documento repetido entre
particiones de ingesta.

La validación de cada fuente se hace por capas, con severidad creciente.
Un documento CVN pasa tres comprobaciones, cada una más estricta que la
anterior:

\begin{itemize}
  \item \textbf{Estructural.} Repite la validación del formato Open CVN ya
  usada en el aterrizaje del Capítulo~\ref{cap:ingesta}, como defensa
  adicional frente a otro posible productor de bronze en el futuro.
  \item \textbf{Por entidad.} Valida cada entrada contra la definición
  estricta de su propio tipo de entidad, el mismo mecanismo que ya
  demostró ser necesario durante la generación sintética.
  \item \textbf{Reglas de negocio.} Comprueba reglas propias de esta
  etapa, como que todo identificador ORCID declarado tenga un dígito de
  control válido y que las fechas de cada periodo respeten su orden.
\end{itemize}

Un registro de ORCID pasa una validación paralela basada en reglas, con el
dígito de control del identificador y la presencia de nombre y de al
menos una afiliación u obra. Los registros de cualquiera de las dos
fuentes que no superan sus reglas de validación se separan a una tabla de
rechazados con la regla concreta que fallaron,
siguiendo el mismo principio de conservar el rastro ya aplicado al
aterrizaje en bronze, y el trabajo falla si la proporción de rechazados de
una fuente supera un umbral del 5\%.

Los registros válidos se normalizan a una forma común antes de la
resolución de identidad, en dos frentes:

\begin{itemize}
  \item \textbf{Nombres.} Se pliegan sin acentos y en minúsculas, y las
  partículas de apellido como «de» o «del» se descartan al construir la
  clave de bloqueo, para que «de la Torre» quede indexado por «torre» y no
  por la partícula.
  \item \textbf{Organizaciones.} Se reducen a su conjunto de palabras
  significativas, ordenado y sin las de relleno, de forma que el orden en
  que aparecen dentro del nombre deje de importar y pueda calcularse una
  similitud por solapamiento de conjuntos.
\end{itemize}

Sobre el subconjunto real verificado con 10\,000 CVN sintéticos generados,
esta etapa produjo 30\,868 registros de persona válidos, 99\,979
afiliaciones y 543\,879 publicaciones, sin ningún registro rechazado.

\section{Resolución de identidad entre fuentes}
\label{sec:procesamiento_resolucion}

Fusionar el registro de un mismo investigador presente tanto en ORCID como
en un documento CVN es el núcleo del problema de fusión planteado en el
Capítulo~\ref{cap:introduccion}. El estado del arte de la resolución de
entidades ofrece, a grandes rasgos, tres familias de enfoque:

\begin{itemize}
  \sloppy
  \item \textbf{Coincidencia determinista por reglas explícitas.} Un
  identificador único compartido y, en su ausencia, nombre normalizado y
  afiliación compartida.
  \item \textbf{Vinculación probabilística de registros.} Calcula una
  probabilidad de coincidencia combinando varios campos con pesos
  calibrados, en lugar de una decisión binaria. La formalización clásica
  es de Fellegi y Sunter \cite{fellegi_sunter_1969}. Splink
  \cite{splink_docs} es una herramienta que la implementa.
  \item \textbf{Resolución basada en aprendizaje automático o
  representaciones vectoriales.} Aprende una función de similitud a
  partir de ejemplos etiquetados. La biblioteca \texttt{dedupe}
  \cite{dedupe_docs} es un ejemplo disponible.
\end{itemize}

La opción adoptada es la coincidencia determinista, en tres reglas
aplicadas por orden de fuerza de evidencia:

\begin{itemize}
  \item \textbf{Regla R1.} Fusiona todo registro que declare el mismo
  identificador ORCID válido, evidencia fuerte porque desambiguar autores
  de forma inequívoca es el objetivo fundacional de ORCID.
  \item \textbf{Regla R2.} La regla de repliegue para un CVN que no
  declara ese identificador. Exige que la clave de bloqueo del CVN,
  apellido normalizado más inicial del nombre, coincida con la de una
  entidad ya anclada por R1, que el nombre del CVN sea compatible
  permitiendo relajar como mucho una de las dos partes, apellido o nombre,
  nunca las dos a la vez, que alguna organización de esa entidad coincida
  tras la normalización, y que exactamente una entidad cumpla las tres
  condiciones.
  \item \textbf{Regla R3.} Deja como entidad propia todo registro que
  ninguna de las dos anteriores fusiona.
\end{itemize}

La Figura~\ref{fig:tfm_resolucion_flujo} resume el orden de aplicación. El
diseño favorece la precisión sobre la cobertura de forma deliberada. Una
fusión equivocada corrompe todo indicador que dependa de esa fusión más
adelante, mientras que una fusión no realizada solo divide una entidad en
dos.

\begin{figure}[H]
  \centering
  \resizebox{0.92\textwidth}{!}{%
  \begin{tikzpicture}[
    font=\small,
    node distance=9mm and 12mm,
    decision/.style={rectangle, draw=azul, thick, dashed, rounded corners,
      text width=4.6cm, align=center, inner sep=5pt},
    outcome/.style={rectangle, draw=azul, thick, rounded corners,
      fill=gris95, text width=4cm, align=center, inner sep=5pt},
    start/.style={rectangle, draw=gris50, text width=4cm, align=center,
      inner sep=5pt},
    >=Stealth
  ]
    \node[start] (inicio) {Registro sin entidad asignada};
    \node[decision, below=of inicio] (d1) {¿Declara un identificador ORCID válido?};
    \node[outcome, right=of d1] (r1) {R1: fusiona con la entidad de ese identificador};
    \node[decision, below=of d1] (d2) {¿Clave de bloqueo, nombre y organización coinciden con una única entidad anclada por R1?};
    \node[outcome, right=of d2] (r2) {R2: fusiona con esa entidad};
    \node[outcome, below=of d2] (r3) {R3: entidad propia, singleton};

    \draw[->] (inicio) -- (d1);
    \draw[->] (d1) -- node[above]{sí} (r1);
    \draw[->] (d1) -- node[left]{no} (d2);
    \draw[->] (d2) -- node[above]{sí} (r2);
    \draw[->] (d2) -- node[left]{no} (r3);
  \end{tikzpicture}}
  \caption{Orden de aplicación de las reglas de resolución de identidad.}
  \label{fig:tfm_resolucion_flujo}
\end{figure}

Frente a la vinculación probabilística o el aprendizaje automático, un
enfoque determinista con evidencia explícita por tipo de enlace produce
resultados auditables campo a campo, una propiedad valiosa en un trabajo
académico donde cada decisión debe poder justificarse de forma directa. La
alternativa probabilística o de aprendizaje automático queda documentada
como trabajo futuro.

El resultado se midió contra un conjunto de verdad de referencia que el
propio generador sintético produce, no contra una anotación externa, sobre
la ejecución de 10\,000 documentos CVN. La Tabla~\ref{tab:tfm_resolucion_resultados}
resume el resultado de cada regla.

\begin{table}[H]
  \centering
  \small
  \renewcommand{\arraystretch}{1.2}
  \begin{tabular}{
    p{0.06\textwidth}
    >{\raggedright\arraybackslash}p{0.26\textwidth}
    p{0.13\textwidth} p{0.13\textwidth} p{0.12\textwidth} p{0.12\textwidth}}
    \hline
    \textbf{Regla} & \textbf{Evidencia} & \textbf{Evaluables} & \textbf{Correctas} & \textbf{Precisión} & \textbf{Cobertura} \\
    \hline
    R1 & Identificador ORCID compartido & 7\,098 & 7\,098 & 100\% & 100\% \\
    R2 & Nombre y afiliación normalizados & 186 & 138 & 95,8\% & 74,2\% \\
    \hline
  \end{tabular}
  \caption{Resolución de identidad sobre la ejecución de 10\,000 documentos CVN.}
  \label{tab:tfm_resolucion_resultados}
\end{table}

Esa cobertura del 74,2\% está en su techo real y no en un defecto de la
regla. Una cuarta parte de los documentos evaluables no declara ninguna
afiliación, y sin afiliación la regla no tiene evidencia posible por
diseño. Un umbral de precisión del 99\% se había fijado como objetivo
antes de medir, y no se alcanza. Exigir más de una organización compartida
para aceptar un nombre acortado sube la precisión al 98,5\% pero baja la
cobertura al 70,4\%, y ningún ajuste del umbral de similitud entre
organizaciones cambia el resultado. Los seis falsos positivos de la tabla
comparten nombre y una única organización con la persona correcta,
indistinguibles con esta evidencia de una coincidencia real. La cifra que
se retiene, por tanto, es la de la tabla, aceptada como el límite honesto
de este diseño en lugar de perseguirse artificialmente con más reglas.

\section{Cálculo de indicadores y publicación en PostgreSQL}
\label{sec:procesamiento_gold}

El segundo trabajo calcula tres indicadores de investigación sobre silver.
El primero cuenta publicaciones por investigador y año. El segundo
construye la línea temporal de afiliaciones de cada entidad, con su rango
de años de carrera. El tercero detecta pares de colaboración entre
entidades que comparten el identificador de una misma obra. Antes de
calcular el primero, es necesario deduplicar las publicaciones por
entidad, porque una misma obra puede aparecer en varios registros
fusionados en una entidad, un CVN y su registro ORCID enlazado, por
ejemplo. Sobre los datos reales verificados, deduplicar redujo 543\,879
filas de publicación a 516\,396 obras distintas. Sin esa deduplicación, el
indicador habría contado un 5,1\% de publicaciones de más.

El indicador de colaboración exige una comprobación previa, porque el
mismo problema de cobertura de la Sección~\ref{sec:procesamiento_resolucion}
puede producir pares de colaboración falsos. Una persona cuyo CVN no llegó
a fusionarse con su propio registro ORCID reportaría las mismas obras
desde dos entidades distintas, y aparecería colaborando consigo misma.
Sobre los 9\,540 pares de entidades que comparten alguna obra en los datos
reales, solo 25, un 0,26\%, resultaron ser ese artefacto, medido contra el
mismo conjunto de verdad de referencia de la sección anterior. El
indicador se construye igualmente, con una marca por par que indica si
alguna de las dos entidades procede de un CVN, para poder restringir el
análisis a colaboraciones solo entre datos reales de ORCID cuando se desee
esa lectura más estricta.

Gold se publica en PostgreSQL a través de tablas de preproducción y un
intercambio atómico. El trabajo crea cada tabla final vacía con su
definición explícita, carga sus filas mediante una escritura JDBC
ordinaria, y solo entonces sustituye la tabla publicada por la nueva
dentro de una única transacción, de forma que el panel de visualización
nunca ve una tabla a medio cargar. Antes de esa sustitución se reconcilia
el número de filas entre Iceberg y PostgreSQL, y cualquier discrepancia
aborta la publicación sin tocar lo ya publicado. Este mecanismo se
verificó deliberadamente matando el proceso de publicación a mitad de
carga. Las tablas ya publicadas permanecieron intactas, y una repetición
posterior del mismo trabajo limpió los restos de la carga interrumpida y
completó la publicación con normalidad.

\section{Verificación de extremo a extremo}
\label{sec:procesamiento_verificacion}

Siguiendo el mismo patrón de doble comprobación de los capítulos
anteriores, el resultado de cada trabajo se contrastó con un cálculo
independiente. Para gold, ese cálculo fue un script en Python normal,
escrito deliberadamente con estructuras de datos distintas a las de Spark,
que recalculó cada tabla a partir de una exportación de silver. Comparado
fila a fila contra el resultado real del flujo, no encontró ninguna
discrepancia sobre 306\,599 filas publicadas en PostgreSQL. Una segunda
ejecución completa del flujo, sobre el mismo bronze, produjo un contenido
idéntico fila a fila en PostgreSQL, lo que confirma que la reconstrucción
completa en cada ejecución es determinista.

La última comprobación midió cuánto dependen los indicadores de los
errores conocidos de la resolución de identidad. Deshacer las 75 fusiones
de la regla R2 que se apoyan en una única organización compartida, el
subconjunto de evidencia más débil, desplaza cada indicador como mucho un
0,58\%. Los indicadores de gold son, por tanto, insensibles en la práctica
a la precisión del 95,8\% medida para la regla R2, aunque esa precisión no
alcance el objetivo inicial del 99\%.
